\chapter{Consistent Learning and Description Length}
\label{sec:p15-occam-learning}

Fitting the observed examples does not, by itself, establish accuracy
on further examples. A learning guarantee relates the sampling mechanism,
the class of candidate functions and the rule that chooses a function.
This chapter supplies the finite-class lemma and Occam-algorithm theorem
of Blumer, Ehrenfeucht, Haussler and Warmuth
\cite[Lemma~2.1, Definition~2.2 and Theorem~2.3]{BlumerEtAl1987Occam}.
It states both the probability and computational assumptions, with
explicit finite-string counting conventions.

\section{Targets, observations and generalization error}

\begin{hbdefinition}{def:p15-occam-model}{The realizable independent-sample model}
Let \((X,\mathcal A)\) be a nonempty measurable input space,
\(Y\) a nonempty finite label set with its discrete sigma-field, and
\(H\subseteq Y^X\) a nonempty finite or countable class of measurable
functions fixed before the sample. The target
\(f\in H\) is fixed. A distribution \(P\) on \(X\) produces independent
observations \(X_1,\ldots,X_m\), all with law \(P\), where \(m\geq1\).
The learner receives
\(S=((X_1,f(X_1)),\ldots,(X_m,f(X_m)))\).
The sample is ordered and may contain repeated inputs.

A hypothesis \(h\in H\) is consistent with \(S\) when
\(h(X_i)=f(X_i)\) for every \(i\). Its generalization error and
empirical error are
\[
 R_P(h,f)=P\{x:h(x)\ne f(x)\},\qquad
 \widehat R_S(h)=\frac1m\sum_{i=1}^m
                       \mathbf1_{\{h(X_i)\ne f(X_i)\}}.
\]
The version space \(V(S)\) is the set of consistent hypotheses.
It is nonempty because it contains \(f\).
A proper consistent learner returns a member of \(V(S)\).
Here \(\varepsilon\in(0,1)\) is an error tolerance and
\(\delta\in(0,1)\) is a failure-probability tolerance. The first
concerns a fresh input under \(P\); the second concerns the sample
under \(P^{\otimes m}\). A learner's selected hypothesis is assumed
measurable as an \(H\)-valued random variable, with the discrete
sigma-field on \(H\). In the source's encoded computational setting,
\(X\) is countable with the full power-set sigma-field; selection
maps on the resulting countable sample space are automatically measurable.
\index{learning!realizable}\index{generalization error}\index{version space}
\end{hbdefinition}

\begin{hbtheorem}{thm:p15-occam-finite}{proved=annex:p15.occam-finite}{Every consistent member of a fixed finite class}
If \(|H|=r<\infty\), then
\[
 P^{\otimes m}\{\exists h\in V(S):R_P(h,f)>\varepsilon\}
      <r(1-\varepsilon)^m\leq r e^{-m\varepsilon}.
\]
In particular every proper consistent learner has error at most
\(\varepsilon\) with probability at least \(1-\delta\) whenever
\[
 m\geq
 \max\left\{1,\left\lceil
       \frac{\ln r+\ln(1/\delta)}{-\ln(1-\varepsilon)}
                         \right\rceil\right\}.
\]
Replacing the denominator by \(\varepsilon\) gives a sufficient
larger bound. These guarantees are uniform in \(P\) and \(f\in H\).
They establish a statistical bound even if finding a consistent
hypothesis is computationally expensive.
\end{hbtheorem}

\begin{hbtheorem}{thm:p15-occam-threshold-example}{proved=annex:p15.occam-threshold-example}{A complete finite version-space calculation}
Let \(X=\{0,1,2\}\), let \(P\) be uniform, and let
\(H=\{h_0,h_1,h_2,h_3\}\), where \(h_t(x)=\mathbf1_{\{x\geq t\}}\).
Take target \(f=h_2\) and tolerance \(\varepsilon=1/4\).
The exact probability that some bad hypothesis remains consistent is
\[
                 2(2/3)^m-(1/3)^m.
\]
The smallest sample size for which this probability is at most
\(1/20\) is \(m=10\): its values at \(m=9,10\) are respectively
\(1023/19683\) and \(2047/59049\).
The generic bound in \cref{thm:p15-occam-finite} gives the
sufficient sample size
\(\lceil\ln80/\ln(4/3)\rceil=16\).
A sample containing inputs zero and two but not one leaves
both \(h_1\) and \(h_2\) consistent, although \(h_1\) has error \(1/3\).
\end{hbtheorem}

\section{What a short description counts}

\begin{hbdefinition}{def:p15-occam-codes}{A fixed representation system}
Let \(\mathcal C\subseteq\{0,1\}^*\) be a set of binary codes with
a fixed decoding map onto \(H\). A code may name one function;
several codes may name the same function. Its length is its number
of bits, and \(\ell(h)\) is the minimum code length of \(h\).
An empty code has length zero. Decoding conventions are fixed
before observing data.

A code family is prefix-free when no code is a proper prefix
of another distinct code. A prefix-free family containing the
empty code contains no other code. For computational claims, restrict
to countable \(X\) with fixed effective finite encodings of inputs
and labels, and assume that testing
one represented hypothesis on one encoded observation is assumed
to take polynomial time in their combined lengths.
\index{description length!fixed encoding}\index{prefix-free code}
\end{hbdefinition}

\begin{hbtheorem}{thm:p15-occam-code-count}{proved=annex:p15.occam-code-count}{Arbitrary strings and prefix-free descriptions}
For real \(L\geq0\), the number of distinct hypotheses with a code
of length at most \(L\) is at most
\[
                 2^{\lfloor L\rfloor+1}-1<2^{L+1}.
\]
For a prefix-free code family, the sharper bound \(2^{\lfloor L\rfloor}\)
holds. Furthermore a prefix-free family satisfies the Kraft inequality
\[
                       \sum_{c\in\mathcal C}2^{-|c|}\leq1.
\]
Choosing one shortest code per hypothesis therefore gives
\(\sum_{h\in H}2^{-\ell(h)}\leq1\).
Without the prefix-free condition, lengths at most one already
allow the three codes consisting of the empty word, zero and one.
\end{hbtheorem}

\section{The Occam algorithm and its full sample bound}

\begin{hbdefinition}{def:p15-occam-algorithm}{The source's polynomial compression requirement}
Fix constants \(c\geq1\) and \(0\leq a<1\).
An Occam algorithm is a proper consistent learner such that,
for every integer target-size bound \(n\geq1\), every \(f\in H\)
with \(\ell(f)\leq n\), and every nonempty sample of \(f\) of
size \(m\), it outputs a code of length at most \(n^c m^a\).
It runs in time polynomial in the total encoded sample length.
The constants and representation system do not depend on
\(f,P,n,m,\varepsilon,\delta\).

The source's polynomial learnability requires a sufficient sample
size polynomial in \(n,1/\varepsilon,1/\delta\), uniform in
\(P\) and \(f\) of size at most \(n\), and polynomial running
time in the encoded sample length. The statistical theorem
below supplies such a sample size. It does not assert that
an efficient Occam algorithm exists for every representation class.
\end{hbdefinition}

\begin{hbtheorem}{thm:p15-occam-compression}{proved=annex:p15.occam-compression}{The complete Occam implication with explicit code conventions}
Set \(q=-\ln(1-\varepsilon)>0\).
For arbitrary binary codes, the Occam algorithm has error at most
\(\varepsilon\) with probability at least \(1-\delta\) if
\[
 m\geq\left\lceil\max\left\{
  1,\ \frac{2\ln(2/\delta)}q,\
       \left(\frac{2n^c\ln2}{q}\right)^{1/(1-a)}
                          \right\}\right\rceil .
\]
For prefix-free codes, replace \(\ln(2/\delta)\) by
\(\ln(1/\delta)\). This gives the source's displayed
constant bound under that counting convention.
For either convention, a sufficient sample size is
\[
 O\!\left(\frac{\ln(1/\delta)}{\varepsilon}
               +\left(\frac{n^c}{\varepsilon}\right)^{1/(1-a)}
               +1\right)
\]
for \(0<\varepsilon,\delta<1\) and \(n\geq1\); the constants depend
only on \(c,a\). Thus the source's polynomial learnability
conclusion follows.
\end{hbtheorem}

\begin{hbtheorem}{thm:p15-occam-minimum-approximation}{proved=annex:p15.occam-minimum-approximation}{The sample-independent case and approximate minimization}
If the output length is at most \(n^c\), independent of \(m\),
then arbitrary codes require only
\[
 m\geq\max\left\{1,\left\lceil
       \frac{(n^c+1)\ln2+\ln(1/\delta)}{-\ln(1-\varepsilon)}
                                      \right\rceil\right\}.
\]
For prefix-free codes, \(n^c+1\) can be replaced by \(n^c\).

Let \(b(S)\) be the minimum length of any consistent hypothesis,
and use \(b_+(S)=\max\{1,b(S)\}\) to handle the empty code.
If a polynomial-time proper consistent algorithm returns length
at most \(b_+(S)^c\), it meets this sample-independent bound
for every target of length at most \(n\geq1\).
Exact minimum-length selection is sufficient when computationally
available, but exact minimization is not required by the theorem.
\end{hbtheorem}

\section{A simultaneous guarantee for a countable class}

\begin{hbtheorem}{thm:p15-occam-weighted}{proved=annex:p15.occam-weighted}{Fixed weights and a guarantee for the selected hypothesis}
Assign weights \(w_h>0\) to a countable \(H\) before sampling,
with \(\sum_h w_h\leq1\). With probability at least \(1-\delta\),
every \(h\in V(S)\) satisfies
\[
 R_P(h,f)\leq
      \frac{\ln(1/w_h)+\ln(1/\delta)}m.
\]
A bound exceeding one is valid but uninformative.
For prefix-free descriptions one may take \(w_h=2^{-\ell(h)}\).
For arbitrary binary descriptions, choosing one shortest code
per hypothesis permits \(w_h=2^{-2\ell(h)-1}\).
These give respectively the numerators
\(\ell(h)\ln2+\ln(1/\delta)\) and
\((2\ell(h)+1)\ln2+\ln(1/\delta)\).
The weights allocate a probability bound over hypotheses; this
theorem does not construct a Bayesian posterior over \(H\).
\end{hbtheorem}

\section{Why the assumptions matter}

\begin{hbtheorem}{thm:p15-occam-boundaries}{proved=annex:p15.occam-boundaries}{Dependence, distribution change and a data-dependent class}
The preceding guarantees cannot in general discard their sampling
and fixed-class assumptions:
\begin{enumerate}
\item On two equally likely inputs, a hypothesis wrong at one
      input remains consistent with probability \(1/2\) if all
      \(m\) observations repeat one random input. The independent
      value would be \(2^{-m}\).
\item A hypothesis can have training error zero and future error
      one when the training and future input distributions differ.
\item After observing an all-zero target on a finite \(N\)-point
      domain, form \(h_S\) that is zero at observed points and
      one elsewhere. The data-dependent class \(\{f,h_S\}\)
      has at most two members and contains the target, but
      returning \(h_S\) has uniform-distribution error at least
      \(1-m/N\). If \(N\geq2m\), that error is at least \(1/2\)
      for every sample.
\end{enumerate}
Thus choosing a short code or small class after observing data
does not establish the fixed-representation premise. Conflicting
labels for the same input can also make the version space empty;
the noiseless realizability guarantee then does not apply.
\end{hbtheorem}

\begin{hbtheorem}{thm:p15-occam-unseen-labels}{proved=annex:p15.occam-unseen-labels}{A finite lower bound without restrictions on the target}
On a domain of \(N\geq1\) points with uniform input distribution,
let the target range over all binary functions. For any learner
using \(m\) labelled observations, and with any independent internal
randomization, some fixed target has expected generalization error
at least
\[
                           \frac12\left(1-\frac mN\right).
\]
When \(N\geq2m\), this is at least \(1/4\).
This lower bound is an editorial comparison: it does not
contradict learnability with sample size depending on the
target representation length.
\end{hbtheorem}

\section{Using a learning guarantee}

\begin{hbmethod}{def:p15-occam-method}{From a representation class to a defensible error claim}
Specify the input population and target labels. Choose the hypothesis
class and encoding before observing the sample. State whether the
learner is proper, returns a consistent hypothesis and has the required
output-size and runtime bounds. Select \(\varepsilon,\delta\), then
choose the sample size from the applicable theorem.

The resulting conclusion concerns error on a fresh input from the
same distribution, with the stated probability over the sample.
It is not exact identification of the target, a guarantee for
adversarial distribution changes, or a philosophical claim that
the shortest explanation must be true.
\end{hbmethod}

The original paper separates its proved implication from its
Section~3 questions about efficient Occam algorithms and a converse.
Its discussion of hard minimum-consistency problems refers to
other works; those reductions and later converse theorems are
separate presentations. Noise, improper learning, agnostic loss
minimization, infinite-dimensional complexity and modern generalization
methods require their own theorems. Bayesian updating under a
specified prior is treated in
\cref{sec:p15-bayesian-updating-prediction}; here the distribution
over training inputs remains arbitrary within the stated model.
